\chapter{Diagramas de estado}
\label{ch:estados}

Los diagramas de estado describen el ciclo de vida de los recursos con estados propios. Las fuentes editables están en \path{docs/diagramas/estado}.

\section{Ciclo de vida de la receta}

\begin{figure}[H]
  \centering
  \includegraphics[width=\textwidth]{build/estado/estado-receta.pdf}
  \caption{Estados de una receta: visibilidad, reporte y eliminación.}
  \label{fig:estado-receta}
\end{figure}

La receta nace privada. La visibilidad y el reporte son atributos del documento, y el borrado es físico: no hay borrado lógico ni cascada.

\begin{itemize}[leftmargin=*]
  \item \texttt{reportado} es un booleano independiente; una receta reportada no cambia de visibilidad ni deja de ser consultable.
  \item No existe estado de moderación ni cuenta administradora.
  \item Al eliminar, se borra la receta y su imagen GridFS, pero quedan huérfanos los comentarios, reseñas y actividades asociados.
\end{itemize}

\section{Ciclo de vida de la imagen}

\begin{figure}[H]
  \centering
  \includegraphics[width=\textwidth]{build/estado/estado-imagen.pdf}
  \caption{Estados de la imagen de una receta.}
  \label{fig:estado-imagen}
\end{figure}

Solo se conserva una imagen vigente por receta. El archivo anterior se elimina de GridFS al reemplazar, quitar o borrar la receta.

\begin{itemize}[leftmargin=*]
  \item La lectura \path{GET /api/recipes/image/:fileId} es pública y usa caché inmutable.
  \item La validación real de bytes (magic bytes) está pendiente; solo se comprueba el tipo MIME declarado y un máximo de 5 MB en el servidor.
\end{itemize}

\section{Estado de la campana de notificaciones}

\begin{figure}[H]
  \centering
  \includegraphics[width=\textwidth]{build/estado/estado-campana-notificaciones.pdf}
  \caption{Estado provisional de la campana de notificaciones.}
  \label{fig:estado-campana}
\end{figure}

La campana es provisional: reutiliza el feed y no persiste el estado de lectura.

\begin{itemize}[leftmargin=*]
  \item El contador de no leídas es la longitud de la lista, no un estado real.
  \item \texttt{markAsRead()} solo modifica el arreglo en memoria del navegador.
  \item El modelo \texttt{Notification} con \texttt{destinatario}, \texttt{leida} e índice \texttt{\{ destinatario: 1, leida: 1, createdAt: -1 \}} está especificado pero no implementado (Fase 7).
\end{itemize}
